
表\ref{tab:sim_metrics}给出了五组算例的定量指标。其中 \(\mathrm{RMS}(e)\) 为所有关节、所有采样时刻位置误差的均方根值。无扰动时，PD 与 PDF 均可实现高精度跟踪，PDF 在终端误差上略低于 PD；在有界扰动下，未启用扰动观测器的 PD 和 PDF 控制均只能实现小邻域跟踪。加入 PTDO 后，PDF+PTDO 算例的终端位置误差降至 \(3.139\times10^{-5}\,\mathrm{rad}\)，终端速度误差降至 \(4.226\times10^{-4}\,\mathrm{rad/s}\)，且最大控制力矩未出现明显增加。

\begin{table}[H]
\centering
\caption{不同控制算例的仿真指标}
\label{tab:sim_metrics}
\begin{tabular}{ccccc}
\toprule
控制器与工况 & \(\|e(10)\|\) & \(\|\dot e(10)\|\) & \(\mathrm{RMS}(e)\) & \(\max|\tau_i|\)\\
 & rad & rad/s & rad & N\,m\\
\midrule
PD，无扰动 & \(4.028\times10^{-9}\) & \(1.802\times10^{-8}\) & \(8.622\times10^{-3}\) & 9.129\\
PDF，无扰动 & \(3.290\times10^{-10}\) & \(1.145\times10^{-9}\) & \(8.622\times10^{-3}\) & 9.129\\
PD，有扰动 & \(1.221\times10^{-2}\) & \(1.197\times10^{-2}\) & \(9.455\times10^{-3}\) & 9.129\\
PDF，有扰动 & \(1.235\times10^{-2}\) & \(1.177\times10^{-2}\) & \(9.456\times10^{-3}\) & 9.129\\
PDF+PTDO，有扰动 & \(3.139\times10^{-5}\) & \(4.226\times10^{-4}\) & \(8.623\times10^{-3}\) & 9.129\\
\bottomrule
\end{tabular}
\end{table}